\documentclass[11pt,a4paper]{article}

\usepackage[margin=2.5cm]{geometry}
\usepackage[T1]{fontenc}
\usepackage{lmodern}
\usepackage{amsmath}
\usepackage{booktabs}
\usepackage{tabularx}
\usepackage{enumitem}
\usepackage{xcolor}
\usepackage{listings}
\usepackage{hyperref}

\hypersetup{colorlinks=true, linkcolor=blue!50!black, urlcolor=blue!50!black}
\lstset{basicstyle=\ttfamily\small, columns=fullflexible, breaklines=true, frame=single,
  framerule=0.3pt, rulecolor=\color{black!30}, backgroundcolor=\color{black!3}, showstringspaces=false}
\newcommand{\code}[1]{\texttt{#1}}

\title{CPU--GPU Co-Scheduling Module for CloudSim 7G\\[4pt]
\large Phase 2 Report: Vertical Slice, the First End-to-End Run}
\author{Chirag \and Hamiz}
\date{22 November 2026}

\begin{document}
\maketitle

\begin{abstract}
Phase~2 connects every layer for the first time. An experiment YAML names a cluster, a workload and
candidate schedulers; jobs are read and labelled, flow through a scheduling broker that builds a
round plan with pluggable policies, run on the GPU layer at speeds from the throughput profiles, and
come out as a per-job result file with a run manifest. The first policies are FIFO ordering, the
fixed CPU:GPU share, consolidated placement with STRICT locality, a non-preemptive lifecycle, and
periodic or event-driven rounds. Hand-computed job completion times are matched, results are
identical across runs, the output is read directly by \code{tools/compare\_runs.py}, and the
architecture's dependency rules are enforced by tests. All 242 tests of the project pass.
\end{abstract}

\tableofcontents

%======================================================================
\section{What runs end to end}
%======================================================================

\begin{enumerate}
  \item \code{ExperimentRunner} reads the YAML (\code{docs/formats.md}) and, for every candidate and
        seed, starts a fresh simulation.
  \item The candidate's plug-ins are built by name from \code{StandardPlugins}; every parameter is
        range-checked and recorded.
  \item The workload CSV is read and labelled; all needed profiles are checked to exist before the run.
  \item The GPU layer builds one host per server; the broker is created and bound to it.
  \item The simulation runs. At the end the runner checks that every job finished and writes
        \code{jobs.csv} and \code{manifest.json} under \code{<output\_dir>/<candidate>/seed-<seed>/}.
\end{enumerate}

%======================================================================
\section{Policy interfaces}
%======================================================================

Each plug point is a small interface in \code{cosched.policy}. Policies read only the cluster view
(\code{ClusterState}), the round plan and the models, and write only to the round plan
(architecture rule~6).

\begin{center}
\begin{tabularx}{\linewidth}{@{}lX@{}}
\toprule
Interface & Decides \\
\midrule
\code{SchedulingTrigger} & When rounds happen: every \code{round\_s}, and/or on job arrival and end. \\
\code{OrderingPolicy} & Which job is served first. \\
\code{AllocationPolicy} & How many cores and how much memory a job asks for on a server group. \\
\code{PlacementPolicy} & Which servers a job runs on, reading free capacity from the plan. \\
\code{LifecycleController} & Whether running jobs keep their allocation or are re-allocated each round. \\
\bottomrule
\end{tabularx}
\end{center}

\subsection{First implementations}
\begin{itemize}
  \item \textbf{FIFO}: by arrival time, ties by id. Every waiting job is tried in that order, and a job
        that does not fit is skipped for this round, so a smaller later job may use GPUs the head job
        cannot.
  \item \textbf{GPU-proportional allocation} (Synergy's ``+fair''): per GPU, server cores and memory
        divided by server GPUs (3 cores and 62.5\,GB on Synergy's server).
  \item \textbf{Consolidated placement, STRICT}: the fullest server that fits (fewest free GPUs, ties
        by index), never spanning servers. SOFT locality comes with Synergy's policies in Phase~3.
  \item \textbf{Non-preemptive lifecycle}: the plan starts with running jobs as they are; rounds only
        place waiting jobs.
  \item \textbf{Triggers}: periodic (default 300\,s, Synergy's \code{round\_duration}) or event-driven.
\end{itemize}

\code{StandardPlugins} is the only class that names concrete policies (rule~7). Each entry states
its parameters' defaults, ranges and sources.

%======================================================================
\section{The broker}
%======================================================================

\code{CoSchedulingBroker} is a CloudSim entity. It extends \code{SimEntity} rather than
\code{DatacenterBroker} as first planned: \code{DatacenterBroker} processes guest-creation
acknowledgements and cloudlet returns itself, and in this design those belong to the GPU layer's
executor. Every event goes to the executor first.

\paragraph{Arrivals.} Jobs arrive in order from the workload; each arrival schedules the next, so the
event queue stays small. Work is the job's iterations; a trace job given as a duration gets one work
unit per second at the fixed share. A job needing more GPUs than the placement can ever give (with
STRICT locality: more than one server's GPUs) stops the run before it starts, with advice to filter
such jobs.

\paragraph{A round.}
\begin{enumerate}
  \item Take a read-only snapshot: time, servers, free resources, active jobs, organisation.
  \item Create the round plan as the lifecycle asks (empty, or carrying running jobs).
  \item Order the jobs not yet in the plan; for each, ask allocation for its demand on each server
        group and placement for servers, and write the result into the plan. A placement that does
        not fit its own plan is a bug and stops the run.
  \item Diff the plan against what is running: \code{START}, \code{RESIZE}, \code{STOP}, or nothing.
  \item For each change compute the job's rate and GPU busy fraction (\code{RateCalculator}) and send
        the commands to the GPU layer in one call.
\end{enumerate}

\paragraph{Ending.} With periodic rounds, an idle cluster skips ahead to the first round at or after
the next arrival; when nothing is running and nothing is left to arrive, the broker schedules no more
rounds and the run ends. The broker never ends the simulation itself (GPU-layer contract rule~6).

\paragraph{Admission control} is a named method that accepts every job, a seam rather than a plug
point, since no paper or trace used so far needs it.

%======================================================================
\section{Outputs}
%======================================================================

\paragraph{\code{jobs.csv}.} One row per job, sorted by id:
\code{job\_id, model, gpus, arrival\_s, start\_s, finish\_s, jct\_s, queue\_s, starts, cpus\_per\_gpu,
mem\_gb\_per\_gpu}. JCT is finish minus arrival, queueing included, as Synergy defines it. Values are
rounded to the microsecond, far below CloudSim's 0.02\,s resolution.

\paragraph{\code{manifest.json}.} Experiment, candidate and seed; git commit (marked \code{+dirty}
with uncommitted changes); SHA-256 of the YAML and the job CSV; model settings; every plug-in with its
resolved parameters; the cluster; the assumptions in words; model warnings; the number of rounds and
commands; and the summary (mean, p99, p99.9 JCT in hours, mean queueing delay, makespan).

%======================================================================
\section{Smoke run}
%======================================================================

\code{presets/smoke/} holds 20 jobs on two Synergy-shaped servers (8 V100, 24 CPUs, 500\,GB), with
real models and iteration times from the converted profiles, and two candidates: FIFO with the fixed
share under 300\,s rounds, and the same with event-driven rounds.

\begin{center}
\begin{tabular}{@{}lrrrr@{}}
\toprule
Candidate & Mean JCT (h) & p99 JCT (h) & Mean queue (s) & Makespan (s) \\
\midrule
FIFO, fixed share, 300\,s rounds & 1.824 & 4.814 & 2408 & 21{,}855 \\
FIFO, fixed share, event-driven & 1.758 & 4.653 & 2170 & 21{,}345 \\
\bottomrule
\end{tabular}
\end{center}

Both runs take about one second together. Event-driven rounds do slightly better, as expected,
because jobs do not wait for the next 300\,s boundary. Spot checks: job~0 (AlexNet, 7,692
iterations at 0.6\,s) finishes at exactly 4615.2\,s; job~1 arrives at 493.2\,s and, under 300\,s
rounds, starts at the 600\,s round. Running the preset twice gives byte-identical \code{jobs.csv}.
\code{tools/compare\_runs.py} reads the output directly; comparing the two candidates reports a
relative difference of the mean of 3.6\%.

%======================================================================
\section{Tests}
%======================================================================

\begin{center}
\begin{tabularx}{\linewidth}{@{}lX@{}}
\toprule
Test & Checks \\
\midrule
event-driven & One 8-GPU server. Job~1 (8 GPUs, 100 iterations at 1/s) runs 0--100\,s; jobs~2 and 3
  (4 GPUs each) wait and both start at 100\,s, finishing at 150\,s and 130\,s. Fixed share is 4 GPUs,
  12 cores, 250\,GB. \\
periodic & Same jobs with 300\,s rounds: jobs~2 and 3 start at 300\,s and finish at 350\,s and 330\,s. \\
FIFO skipping & A later 4-GPU job runs beside a running job while the 8-GPU head job waits. \\
placement & A 2-GPU job joins the half-full server; the 8-GPU job gets the empty one. \\
too large & A 16-GPU job on 8-GPU servers stops the run before it starts. \\
plan diff & Start, resize, stop, and no change for unchanged jobs. \\
smoke run & Both candidates finish all 20 jobs; files and manifest written; resolved parameters
  recorded; event-driven mean JCT is not worse than periodic. \\
repeatability & Two runs give identical \code{jobs.csv}. \\
bad config & Unknown policy name (error lists known names) and out-of-range \code{round\_s}. \\
architecture (5) & Policies depend only on core, models and policy interfaces; core depends on nothing
  else in the module; models do not depend on policies or the broker; no policy implementation uses
  another; only \code{StandardPlugins} names concrete policies. \\
\bottomrule
\end{tabularx}
\end{center}

\begin{center}
\begin{tabular}{@{}lr@{}}
\toprule
Module & Tests (all passing) \\
\midrule
\code{cloudsim} (unchanged) & 167 \\
\code{cloudsim-examples} & 13 \\
\code{cloudsim-gpu} & 16 \\
\code{cloudsim-cosched} & 46 \\
\midrule
Total & 242 \\
\bottomrule
\end{tabular}
\end{center}

%======================================================================
\section{Next step}
%======================================================================
Phase~3 adds Synergy's own policies: Greedy and Tune allocation (Tune as a joint allocation and
placement policy that may shrink co-located jobs through the round plan), the round-based lifecycle,
SOFT locality, and the SRTF, LAS and FTF orderings, with a \code{synergy-replica} preset. Phase~4 then
replays Synergy's reference run and must match its mean JCT within 10\%.

\end{document}
